This review follows a structured, transparent protocol for a \emph{methods-to-requirements}
mapping. We frame it as a \emph{structured narrative (scoping) review} with a documented,
reproducible search rather than a PRISMA systematic review: the central artifact is an
interpretive crosswalk we make auditable, not a meta-analysis or exhaustive field
enumeration. The protocol was fixed before corpus collection; the search, corpus, coding
ledger, and matrix are released so the synthesis is inspectable.

\subsection{Review Questions}
\textbf{RQ1.} What technical AI evaluation and assessment methods exist, organized by what
they measure? \textbf{RQ2.} Which technical evaluation method could \emph{support an
evidentiary case for conformity} with which requirement of (a)~the NIST AI RMF~1.0,
(b)~the EU AI Act high-risk obligations (Articles~9--15), and (c)~ISO/IEC~42001:2023, and
with what strength of \emph{technical relevance}? \textbf{RQ3.} Which requirements are
weakly served or unserved by technical evaluation methods---i.e., where is the
\emph{measurement gap}?

\subsection{Scope}
We include high-risk and high-stakes AI/ML systems, including large language models (LLMs)
and agentic systems; technical evaluation methods that yield artifacts relevant to
conformity; and the three named frameworks with their official companions. We exclude pure
legal or policy commentary lacking a technical method, consumer/low-risk systems, and
sector regulation outside the three frameworks (cited only for context).

\subsection{Information Sources and Search}
We searched IEEE~Xplore, the ACM Digital Library, Scopus, Web~of~Science, and arXiv
(cs.LG, cs.AI, cs.CY, cs.CR, stat.ML), with forward/backward snowballing via
Google~Scholar. Regulatory requirements were taken from the official corpus: the
NIST~AI~RMF~1.0 (NIST~AI~100-1)~\cite{nist2023airmf} and its Generative~AI
Profile~\cite{nist2024genai}; Regulation~(EU)~2024/1689~\cite{euaiact2024}; and
ISO/IEC~42001:2023~\cite{iso42001} with companion standards. The Boolean search template
combined system terms (AI, machine learning, LLM, foundation model, AI agent) with
evaluation terms (evaluation, assessment, audit, benchmark, testing, assurance, conformity
assessment) and method terms (fairness, robustness, safety, red-teaming, interpretability,
calibration, data governance, model card, human oversight, monitoring), plus a dedicated
governance string. The search window was 2015--2026 with an LLM/agentic emphasis on
2022--2026; the language was English; the last search was conducted in June~2026.
Table~\ref{tab:search} lists the concept groups, and the per-database Boolean query
strings---with the term groups defined and substituted for each database---and the search
dates are released for reproducibility in the companion artifact.

\subsection{The Review Corpus and Its Reconciliation}
To keep the corpus auditable we distinguish four record pools. (i)~A \emph{structured
search} across the twelve families returned candidate technical-method records; after
title/abstract and full-text screening and an identifier/citation-verification step
(below), this yielded \textbf{141} verified technical-method references across families
M1, M3--M11, and governance/positioning; privacy-evaluation methods (M12) were verified
separately as a set of \textbf{6} references, so the released method index contains
\textbf{147} records in total. (ii)~For the
fairness family (M2) we \emph{reuse} a previously published, independently verified
\textbf{202}-entry agent-fairness corpus released separately by the authors; because this
is larger than the bespoke per-family searches, we treat it as a reused supplement and
report a sensitivity analysis with and without it. (iii)~\textbf{Six} official
regulatory/standards texts (the NIST, EU, and ISO anchors and companions) were
hand-collected. (iv)~Two unpublished author manuscripts are mentioned in-text only as
non-independent context and are not listed in the reference list. The complete verified corpus ($\approx$355 records: 147 method $+$ 202 reused $+$ 6
official) is released as Supplementary
Information and in the companion repository, each record carrying metadata, source status,
and family assignment, so that every included record is inspectable. The article's
\emph{reference list} cites approximately \textbf{180} records discussed in the narrative; the
remaining reused-corpus entries are released in full rather than enumerated inline.
Figure~\ref{fig:prisma} shows the record flow, every transition of which is reconciled in
the released ledger.

\emph{Family-level source counts} (newly verified method references) are M1~16, M3~14,
M4~10, M5~13, M6~11, M7~15, M8~12, M9~14, M10~14, M11~12, governance/positioning~10, and
privacy (M12)~6; the fairness family (M2) draws on the reused 202-entry corpus. \emph{Sensitivity:} because the matrix codes
\emph{relevance type} rather than counts, removing the reused fairness corpus changes no
matrix cell and no coverage-breadth rating; its only effect is the relative prominence of
action-level fairness in the research agenda (Section~\ref{sec:gaps}), which we flag.

\subsection{Selection Criteria}
A record was \emph{included} if it (i)~proposed, surveyed, or empirically evaluated a
technical AI evaluation method, or was an official regulatory/standards text;
(ii)~addressed at least one of the twelve method families or framework requirements;
(iii)~had retrievable full text; and (iv)~was in English. A record was \emph{excluded} if
it was opinion/marketing without a method, purely legal analysis without technical
evaluation content, a duplicate, not retrievable, or off-scope (Table~\ref{tab:incl}).

\subsection{Data Extraction}
For each included item we extracted a structured record: bibliographic metadata (authors,
year, venue, DOI/arXiv identifier), a one-line statement of the method contributed, the
assigned method family ($M1$--$M12$), the candidate mapped requirement(s) per framework,
and a technical-relevance code.

\subsection{Coding Protocol and Reliability}
Mapping proceeded in three passes. (1)~\emph{Family assignment}: each method is assigned a
primary family and optional secondary families. (2)~\emph{Requirement mapping}: for every
(method, requirement) pair we assign a technical-relevance level grounded in the official
clause text. \emph{Primary} (\direct): the method produces an artifact that the requirement
chiefly asks for---e.g., an adversarial-robustness report for the EU AI Act Article~15
robustness facet. \emph{Supporting} (\partialev): the artifact feeds a required outcome or
document but is insufficient alone---e.g., a benchmark result is an input to the Article~9
risk-management system but does not constitute it. \emph{Contextual} (\indirect): the
artifact informs the requirement without serving as its evidence---e.g., interpretability
analysis informs governance documentation. A blank denotes non-applicability. The
supporting-versus-blank boundary for process obligations is decided by whether the method's
output is a named input to the obligation's process (supporting) or not (blank).
(3)~\emph{Structured re-examination}: a second pass challenged every non-blank cell against
the clause text and downgraded over-claims (e.g., performance$\rightarrow$Article~11,
safety$\rightarrow$Article~9, the technical$\rightarrow$GOVERN cells, and the privacy
recoding); the full cell-level change log is released.

\emph{Reliability.} A second author independently coded a stratified sample of 40 of the
199 non-blank cells ($\approx$20\%), blind to the first-pass codes. Raw exact agreement was
52.5\% (Cohen's $\kappa=0.18$). We report the unweighted $\kappa$ transparently but note it
is deflated here by two well-understood effects: the marginal distribution is heavily skewed
toward \emph{supporting} (the ``$\kappa$ paradox'', under which high agreement on an
unbalanced category still yields a low $\kappa$), and the scale is a fine three-level ordinal
one on which exact-match $\kappa$ penalizes adjacent codes as harshly as opposite ones. Because
the codes are ordinal, the chance-corrected quadratically-weighted $\kappa$ is the appropriate
measure: all 19 disagreements were a single adjacency level apart, with \emph{no}
primary$\leftrightarrow$contextual conflicts, yielding a quadratically-weighted
$\kappa=0.44$ (moderate). The scheme therefore separates strong from
weak relevance reliably, while the supporting-versus-adjacent boundary is interpretive.
Disagreements were adjudicated to a consensus code; the second
coder's systematically more conservative treatment of governance cells was adopted across
the GOVERN column.

\emph{External validation.} To obtain a check independent of the author team, an external
reviewer (not an author of this study; acknowledged below), blind to the first-pass codes,
recoded the same 40-cell sample. Agreement with the first-pass codes was 57.5\% exact, with
Cohen's $\kappa=0.23$ and a quadratically-weighted $\kappa=0.49$ (computed on the same
relevance scale, here including a single not-applicable code); as in the in-team check, all
17 disagreements were a single adjacency level apart, with \emph{no}
primary$\leftrightarrow$contextual conflicts. The independent recoding
reproduces the same pattern---modest exact agreement on the fine ordinal scale, reliable
separation of strong from weak relevance---which supports presenting the crosswalk as a
\emph{proposed}, externally-checked artifact; broader multi-expert validation would
strengthen it further (Section~\ref{sec:limits}). Both the second-coder and external-validator
cell-level ledgers (\texttt{coding\_audit.csv}, \texttt{external\_validation.csv}) are
released as Supplementary Information.

\subsection{Coverage Breadth}
Per requirement we derive a \emph{coverage-breadth} indicator strictly from the number of
method families with \emph{primary} relevance: \emph{broad} ($\geq$3 families),
\emph{moderate} (1--2), and \emph{narrow} (0). We emphasize that this measures breadth of
mapped families, \emph{not} the methodological maturity, validation quality, or deployment
evidence of the underlying methods; it feeds the gap analysis (Section~\ref{sec:gaps}).

\subsection{Standards Access and Citation Integrity}
Every regulatory obligation is cited to its official source with the specific article,
subcategory, or control identifier. For ISO/IEC~42001:2023 we mapped to the published
Annex~A control structure (numbers and titles) and corroborating secondary analyses of the
standard; protected text is not reproduced beyond permissible quotation. No reference
entered the bibliography without a verified DOI, arXiv identifier, or official document URL
confirmed to resolve; reused entries inherited prior verification and were spot-checked.

\subsection{Threats to Validity}
The crosswalk is interpretive: the frameworks are outcome- and process-oriented and do not
prescribe specific methods, so our cells are reasoned engineering interpretations of
technical relevance, not legal compliance determinations, and a different coder could
assign some adjacent codes differently (the reliability analysis quantifies this). The
frameworks also evolve: EU AI Act application dates were revised by a May~2026 political
agreement, delegated acts and harmonised standards remain in progress, and ISO/IEC
companion standards are emerging (Section~\ref{sec:background}); this review therefore
reflects the state of affairs as of June~2026. Database coverage, English-only inclusion,
preprint inclusion, and the reused fairness corpus introduce selection effects, partly
addressed by the sensitivity analysis above. This article is a technical review and does
\emph{not} constitute legal advice.
